\section{The quantum seed realization}\label{core:section:001}

\paragraph{Conventions.}
Unless the zero space is explicitly included as a boundary case, quantum state spaces are nonzero complex Hilbert spaces of finite dimension. A complete packet family is a finite collection satisfying
\[
P_\alpha^*=P_\alpha,\qquad
P_\alpha P_\beta=\delta_{\alpha\beta}P_\alpha,\qquad
\sum_\alpha P_\alpha=I.
\]
All operator norms in the analytic estimates are induced by the norm on the Hilbert space. The estimates do not use the entrywise maximum norm or the Frobenius norm. We write $[A,B]=AB-BA$ and use units with $\hbar=1$. Parameter derivatives are real derivatives. A matrix $X$ is supported on a code projector $P$ when $PXP=X$.

The common operator sequence of the series becomes
\[
\boxed{
\begin{gathered}
\text{quantum seed}
\longrightarrow
\text{gate or observable address}
\longrightarrow
\text{symmetry projection}\\
\longrightarrow
\text{packet transfer}
\longrightarrow
\text{constraint}
\longrightarrow
\text{refinement}
\longrightarrow
\text{reconstruction}
\end{gathered}
}
\]

\begin{figure}[htbp]
\centering
\includegraphics[width=0.98\textwidth]{01_quantum_seed_architecture.pdf}
\caption{The diagram shows the quantum realization of the common sequence of operators in the seed calculus. The same finite packet language describes states, gates, Hamiltonians, channels, geometric projectors, and effective dynamics.}
\label{core:fig:architecture}
\end{figure}

\begin{definition}[Finite quantum seed system]\label{core:definition:001}
A finite quantum seed system is a tuple
\[
\mathfrak Q=
\left(
\cH,
\psi_0,
\cO,
\cR,
\{P_{\ell,\alpha}\},
H,
\tau
\right),
\]
where $\cH$ is a finite-dimensional Hilbert space, $\psi_0\in\cH$ is a nonzero seed, $\cO$ is the declared operator grammar, $\cR$ is a collection of algebraic or physical constraints, $\{P_{\ell,\alpha}\}$ is an exact packet system at each refinement level, $H=H^*$ is a Hamiltonian when dynamics are present, and $\tau$ specifies the topology used in a later completion.
\end{definition}

The generated state space is
\[
\cH(\psi_0;\cO)
=
\operatorname{span}
\left\{
O_r\cdots O_1\psi_0:
O_j\in\cO,
\ r\ge0
\right\}.
\]
The quantum seed is cyclic for the grammar when
\[
\cH(\psi_0;\cO)=\cH.
\]
This is a statement about the pair $(\psi_0,\cO)$, not about the vector alone. A state may be cyclic for the full matrix algebra and noncyclic for a restricted physical gate set.

\section{Quantum seed rank and cyclicity}\label{core:section:002}

\subsection{Finite observable algebras}

Let a finite-dimensional unital $*$-subalgebra of operators act in its standard multiplicity form
\[
\cA
\cong
\bigoplus_{i=1}^r
\left(M_{d_i}(\bbC)\otimes I_{m_i}\right)
\]
on
\[
\cH
\cong
\bigoplus_{i=1}^r
\left(\bbC^{d_i}\otimes\bbC^{m_i}\right).
\]
Here every $d_i$ is a positive integer and every multiplicity $m_i$ is a nonnegative integer. An empty block family represents the zero space and has seed rank zero. The maximum over an empty block family is also assigned the value zero. A vector $\psi_i\in\bbC^{d_i}\otimes\bbC^{m_i}$ is identified with a $d_i\times m_i$ coefficient matrix $C_i$.

\begin{theorem}[Orbit dimension of one quantum seed]\label{core:theorem:001}\label{core:thm:seed-orbit}
For integers $d\ge1$ and $m\ge0$, consider the action of $M_d(\bbC)\otimes I_m$ on $\bbC^d\otimes\bbC^m$. Let $C$ be the coefficient matrix of a vector $\psi$. Then
\[
\boxed{
\dim\operatorname{span}\{(A\otimes I_m)\psi:A\in M_d(\bbC)\}
=d\,\rank C.
}
\]
The seed is cyclic if and only if $\rank C=m$, which is possible only when $m\le d$.
\end{theorem}

\begin{proof}\phantomsection\label{core:proof:001}
Under the identification by coefficient matrices, $(A\otimes I_m)\psi$ corresponds to $AC$. Every row of $AC$ lies in the row space of $C$. Conversely, any $d\times m$ matrix whose rows lie in the row space of $C$ equals $AC$ for a suitable $A$: choose each row of $A$ to produce the desired row combination. The orbit therefore consists of all matrices with $d$ independently chosen rows in the $\rank C$-dimensional row space of $C$, and has dimension $d\rank C$. It equals the full $dm$-dimensional space exactly when $\rank C=m$.
\phantomsection\label{core:proof:001:end}\end{proof}

\begin{theorem}[Quantum seed-rank formula]\label{core:theorem:002}\label{core:thm:seed-rank}
For
\[
\cA=\bigoplus_i(M_{d_i}(\bbC)\otimes I_{m_i})
\quad\text{on}\quad
\cH=\bigoplus_i(\bbC^{d_i}\otimes\bbC^{m_i}),
\]
the minimum number of seeds required to generate $\cH$ is
\[
\boxed{
\operatorname{srank}_{\cA}(\cH)
=
\max_i\left\lceil\frac{m_i}{d_i}\right\rceil.
}
\]
\end{theorem}

\begin{proof}\phantomsection\label{core:proof:002}
Each seed contributes a coefficient matrix of rank at most $d_i$ in the $i$th block. Thus $k$ seeds can cover at most a $kd_i$-dimensional subspace of the multiplicity factor $\bbC^{m_i}$, so $k\ge\lceil m_i/d_i\rceil$ for every $i$.

For the converse, take
\[
k=\max_i\left\lceil\frac{m_i}{d_i}\right\rceil.
\]
In each block, partition a basis of $\bbC^{m_i}$ into at most $k$ groups of size at most $d_i$. Choose the coefficient matrix of seed $a$ so that its nonzero rows span the $a$th group. The combined row spaces of the $k$ seeds then span $\bbC^{m_i}$ in every block. The preceding theorem shows that their $\cA$-orbits generate the whole direct sum.
\phantomsection\label{core:proof:002:end}\end{proof}

\begin{corollary}[Local orbit and Schmidt rank]\label{core:corollary:001}\label{core:thm:schmidt}
Let $\psi\in\cH_A\otimes\cH_B$, with $\dim\cH_A=d_A\ge1$, and let $r$ be its Schmidt rank. The reduced positive operator on $B$ is $\rho_B=\Tr_A(\ket\psi\bra\psi)$. Normalization of $\psi$ does not affect its support. Then
\[
\boxed{
\dim\operatorname{span}\{(A\otimes I_B)\psi:A\in M_{d_A}(\bbC)\}
=d_A r.
}
\]
In particular, local operators on $A$ generate the complete bipartite Hilbert space exactly when the reduced state on $B$ has full support.
\end{corollary}

\begin{proof}\phantomsection\label{core:proof:003}
A Schmidt decomposition of length $r$ gives a coefficient matrix of rank $r$. Applying Theorem~\ref{core:thm:seed-orbit} gives the orbit dimension $d_A r$. The range of $\rho_B$ is the span of the $r$ Schmidt vectors on $B$, so full support is equivalent to $r=\dim\cH_B$. Because $d_A>0$, this is equivalent to the orbit having the full bipartite dimension.
\phantomsection\label{core:proof:003:end}\end{proof}

The conclusion concerns the algebraic span obtained from arbitrary local operators. A normalized unitary gate orbit need not contain every vector, and the span statement does not supply a deterministic protocol for preparing arbitrary remote states. Entanglement controls the multiplicity support appearing in that algebraic span.

\subsection{A computational-basis seed}

\begin{theorem}[One seed for the computational basis]\label{core:theorem:003}\label{core:thm:computational}
For an integer $n\ge0$, let
\[
\psi_0=\ket{0}^{\otimes n}
\]
and let $X_j$ denote the Pauli $X$ operator on qubit $j$. Then
\[
\left\{
X_1^{a_1}\cdots X_n^{a_n}\psi_0:
(a_1,\ldots,a_n)\in\{0,1\}^n
\right\}
\]
is the complete computational basis. Hence the seed is cyclic under the algebra generated by the $X_j$.
\end{theorem}

\begin{proof}\phantomsection\label{core:proof:004}
For a binary address $a=(a_1,\ldots,a_n)$,
\[
X_1^{a_1}\cdots X_n^{a_n}\ket{0\cdots0}
=
\ket{a_1\cdots a_n}.
\]
Distinct binary addresses produce distinct orthonormal basis vectors. There are $2^n$ of them, equal to the dimension of the Hilbert space.
\phantomsection\label{core:proof:004:end}\end{proof}

\begin{certificate}[Three-qubit seed orbit]\label{core:certificate:001}
The original numerical implementation generated all eight computational basis states from $\ket{000}$ using the three local $X$ address operators. It reported an orbit rank of $8/8$.
\end{certificate}

\section{Cyclic symmetry packets}\label{core:section:003}

\subsection{Character projectors}

Let $N\ge1$ be an integer, and let $S$ be a unitary symmetry whose order divides $N$:
\[
S^N=I.
\]
Set $\omega=e^{2\pi i/N}$.

\begin{theorem}[Cyclic quantum projectors]\label{core:theorem:004}\label{core:thm:cyclic}
The operators
\[
\boxed{
P_j
=
\frac1N\sum_{r=0}^{N-1}\omega^{-jr}S^r,
\qquad
0\le j<N,
}
\]
are self-adjoint, pairwise orthogonal projectors satisfying
\[
P_jP_k=\delta_{jk}P_j,
\qquad
\sum_{j=0}^{N-1}P_j=I,
\qquad
SP_j=\omega^jP_j.
\]
\end{theorem}

\begin{proof}\phantomsection\label{core:proof:005}
Using $S^{r+s}=S^{r+s\bmod N}$ and the finite character orthogonality relation,
\[
\frac1N\sum_{r=0}^{N-1}\omega^{(k-j)r}=\delta_{jk},
\]
we obtain $P_jP_k=\delta_{jk}P_j$. Summing over $j$ cancels every term indexed by $r\ne0$ and retains the $r=0$ term, which proves completeness. Unitarity and the change of residue $r\mapsto-r$ give
\[
P_j^*=\frac1N\sum_{r=0}^{N-1}\omega^{jr}S^{-r}=P_j.
\]
Finally, multiplication by $S$ shifts the summation index and contributes the factor $\omega^j$.
\phantomsection\label{core:proof:005:end}\end{proof}

For parity,
\[
S=Z^{\otimes n},
\qquad
P_{\mathrm{even}}=\frac{I+S}{2},
\qquad
P_{\mathrm{odd}}=\frac{I-S}{2}.
\]

\subsection{Tensor refinement}

Fix $N\ge1$. Suppose every factor $a$ has a complete orthogonal family $P^{(a)}_0,\ldots,P^{(a)}_{N-1}$. The fine tensor packets are
\[
P_{\mathbf j}^{\mathrm{fine}}
=
P^{(1)}_{j_1}\otimes\cdots\otimes P^{(m)}_{j_m}.
\]

\begin{theorem}[Tensor/polyphase refinement]\label{core:theorem:005}\label{core:thm:tensor}
For each residue $r\in\bbZ/N\bbZ$, define
\[
\boxed{
P_r^{\mathrm{coarse}}
=
\sum_{j_1+\cdots+j_m\equiv r\pmod N}
P^{(1)}_{j_1}\otimes\cdots\otimes P^{(m)}_{j_m}.
}
\]
Then the $P_r^{\mathrm{coarse}}$ are complete orthogonal projectors, and every coarse packet is the exact sum of its fine descendants.
\end{theorem}

\begin{proof}\phantomsection\label{core:proof:006}
The fine tensor projectors are mutually orthogonal and sum to the identity. Grouping the fine addresses by their total residue partitions the full address set. Sums over disjoint sets of orthogonal projectors remain orthogonal projectors, and the union of the residue classes is the full fine address set.
\phantomsection\label{core:proof:006:end}\end{proof}

\begin{figure}[htbp]
\centering
\includegraphics[width=0.86\textwidth]{10_tensor_parity_refinement.pdf}
\caption{The eight computational packets of three qubits refine the two total-parity packets exactly.}
\label{core:fig:tensor-refinement}
\end{figure}

\begin{certificate}[Exact parity refinement]\label{core:certificate:002}
For three qubits, the local computational projectors were grouped by the parity of their Hamming weight. The resulting coarse projectors agreed with $(I\pm Z^{\otimes3})/2$ with residual $0$.
\end{certificate}

\section{Quantum circuits as packet paths}\label{core:section:004}

Let
\[
U=U_mU_{m-1}\cdots U_1
\]
be a circuit with $m\ge1$, and let $\{P_\alpha\}$ be a complete packet system. For the empty circuit $m=0$, set $U=I$. Its boundary block is $P_\beta P_\alpha=\delta_{\beta\alpha}P_\alpha$.

\begin{theorem}[Exact circuit path expansion]\label{core:theorem:006}\label{core:thm:paths}
For source packet $\alpha_0$ and target packet $\alpha_m$,
\[
\boxed{
P_{\alpha_m}UP_{\alpha_0}
=
\sum_{\alpha_1,\ldots,\alpha_{m-1}}
P_{\alpha_m}U_mP_{\alpha_{m-1}}
\cdots
P_{\alpha_1}U_1P_{\alpha_0}.
}
\]
Summing over the source and target packets reconstructs the complete circuit operator.
\end{theorem}

\begin{proof}\phantomsection\label{core:proof:007}
Insert the identity resolution $I=\sum_\alpha P_\alpha$ between every adjacent pair of gates. Distributing the finite sums gives the path expansion. Summing the left and right boundary projectors reconstructs $U$.
\phantomsection\label{core:proof:007:end}\end{proof}

\begin{theorem}[Packetwise unitarity conservation]\label{core:theorem:007}\label{core:thm:unitarity}
If $U$ is unitary and $U_{\beta\alpha}=P_\beta UP_\alpha$, then
\[
\boxed{
\sum_\beta U_{\beta\alpha}^*U_{\beta\gamma}
=
\delta_{\alpha\gamma}P_\alpha.
}
\]
Similarly,
\[
\sum_\alpha U_{\beta\alpha}U_{\delta\alpha}^*
=
\delta_{\beta\delta}P_\beta.
\]
\end{theorem}

\begin{proof}\phantomsection\label{core:proof:008}
Using completeness and unitarity,
\[
\sum_\beta U_{\beta\alpha}^*U_{\beta\gamma}
=
P_\alpha U^*\left(\sum_\beta P_\beta\right)UP_\gamma
=
P_\alpha P_\gamma.
\]
For the second identity, sum over the source packets and use unitarity in the reverse order.
\phantomsection\label{core:proof:008:end}\end{proof}

\begin{example}[Bell seed circuit]\label{core:example:001}
In this example only, $H$ denotes the Hadamard gate, not a Hamiltonian. Starting from $\ket{00}$,
\[
(H\otimes I)\ket{00}
=
\frac{\ket{00}+\ket{10}}{\sqrt2},
\]
and then
\[
\operatorname{CNOT}(H\otimes I)\ket{00}
=
\frac{\ket{00}+\ket{11}}{\sqrt2}.
\]
The parity packets show that the final state lies entirely in the even sector, while the intermediate circuit path passes through both even and odd addresses.
\end{example}

\begin{figure}[htbp]
\centering
\begin{minipage}[t]{0.48\textwidth}
\centering
\includegraphics[width=\textwidth]{02_bell_seed_probabilities.pdf}
\captionof{figure}{The Bell state generated from the seed $\ket{00}$.}
\label{core:fig:bell}
\end{minipage}\hfill
\begin{minipage}[t]{0.48\textwidth}
\centering
\includegraphics[width=\textwidth]{04_circuit_packet_paths.pdf}
\captionof{figure}{Circuit reconstruction as a sum over packet paths.}
\label{core:fig:circuit-paths}
\end{minipage}
\end{figure}

\begin{certificate}[Bell and path reconstruction]\label{core:certificate:003}
The original numerical implementation reported residuals of $0$ for both the Bell state and its complete reconstruction from parity paths.
\end{certificate}

\section{Symmetry-resolved Hamiltonians and leakage}\label{core:section:005}

Let $H=H^*$ and let $P$ be a packet projector. If $[H,P]=0$, then $P\cH$ is invariant under $e^{-itH}$.

\begin{theorem}[Exact symmetry packets]\label{core:theorem:008}\label{core:thm:symmetry}
If $[H,P_\alpha]=0$ for every packet, then
\[
H=\sum_\alpha P_\alpha HP_\alpha
\]
and
\[
e^{-itH}
=
\sum_\alpha P_\alpha e^{-itH}P_\alpha.
\]
No amplitude transfers between distinct packets.
\end{theorem}

\begin{proof}\phantomsection\label{core:proof:009}
The commutation relation gives $P_\beta HP_\alpha=0$ for $\beta\ne\alpha$. Every polynomial in $H$ has the same block form, and the finite-dimensional exponential is the norm limit of its polynomial series.
\phantomsection\label{core:proof:009:end}\end{proof}

For two qubits and real parameters $J,h$, consider
\[
H
=
JX\otimes X
+h(Z\otimes I+I\otimes Z).
\]
In the ordered even basis $\{\ket{00},\ket{11}\}$ and odd basis $\{\ket{01},\ket{10}\}$,
\[
H_{\mathrm{even}}
=
\begin{pmatrix}
2h&J\\
J&-2h
\end{pmatrix},
\qquad
H_{\mathrm{odd}}
=
\begin{pmatrix}
0&J\\
J&0
\end{pmatrix}.
\]
Therefore
\[
\Spec H_{\mathrm{even}}
=\{\pm\sqrt{4h^2+J^2}\},
\qquad
\Spec H_{\mathrm{odd}}=\{\pm J\}.
\]

\begin{figure}[htbp]
\centering
\includegraphics[width=0.66\textwidth]{03_parity_block_hamiltonian.pdf}
\caption{The Hamiltonian becomes exactly block diagonal when the even and odd basis vectors are grouped.}
\label{core:fig:parity-hamiltonian}
\end{figure}

\subsection{Leakage under broken symmetry}

Set $Q=I-P$.

\begin{theorem}[Commutator leakage bound]\label{core:theorem:009}\label{core:thm:leakage}
For every self-adjoint $H$, orthogonal projector $P$, and real time $t$,
\[
\boxed{
\norm{Qe^{-itH}P}
\le
|t|\,\norm{[H,P]}.
}
\]
In particular, exact commutation gives zero leakage for all time.
\end{theorem}

\begin{proof}\phantomsection\label{core:proof:010}
Set $F(s)=e^{isH}Pe^{-isH}$. Differentiating the exponential gives
\[
F'(s)=i\,e^{isH}[H,P]e^{-isH},\qquad F(0)=P.
\]
Unitary conjugation preserves the operator norm, so $\|F'(s)\|=\|[H,P]\|$. The mean-value estimate on the interval between $0$ and $t$ yields
\[
\|F(t)-P\|\le |t|\,\|[H,P]\|.
\]
Finally, $e^{itH}Qe^{-itH}P=(P-F(t))P$ and $\|P\|\le1$. Left multiplication by $e^{itH}$ preserves the norm and gives the desired inequality. This argument includes positive, negative and zero time.
\phantomsection\label{core:proof:010:end}\end{proof}

\begin{figure}[htbp]
\centering
\includegraphics[width=0.78\textwidth]{05_leakage_bound.pdf}
\caption{A perturbation that breaks parity creates nonzero transfer between the packets. The measured leakage remains below the commutator bound.}
\label{core:fig:leakage}
\end{figure}

\begin{certificate}[Symmetry and negative control]\label{core:certificate:004}
For the Hamiltonian that preserves parity, the norm of the block between distinct packets and the maximum measured leakage were both $0$. Adding an $X\otimes I$ perturbation broke parity and produced a maximum leakage of approximately $0.79196$. The maximum violation of the proved bound remained $0$.
\end{certificate}

\section{Moving packets, connections, and curvature}\label{core:section:006}

\subsection{Moving quantum coordinates}

\begin{theorem}[Moving-frame Hamiltonian]\label{core:theorem:010}\label{core:thm:moving}
Let $U(t)$ and $\psi(t)$ be differentiable on an open real time interval, let $U(t)$ be unitary, and let $H(t)=H(t)^*$. Suppose
\[
i\dot\psi(t)=H(t)\psi(t).
\]
Writing
\[
\psi(t)=U(t)y(t)
\]
gives
\[
\boxed{
i\dot y(t)
=
\left(U(t)^*H(t)U(t)-iU(t)^*\dot U(t)\right)y(t).
}
\]
Both the connection term $-iU^*\dot U$ and the effective Hamiltonian $U^*HU-iU^*\dot U$ are Hermitian.
\end{theorem}

\begin{proof}\phantomsection\label{core:proof:011}
Differentiate $\psi=Uy$:
\[
i\dot U y+iU\dot y=HUy.
\]
Multiplying by $U^*$ gives the formula. Differentiating $U^*U=I$ shows $U^*\dot U$ is skew-Hermitian, so $-iU^*\dot U$ is Hermitian.
\phantomsection\label{core:proof:011:end}\end{proof}

\begin{figure}[htbp]
\centering
\includegraphics[width=0.70\textwidth]{08_moving_frame_connection.pdf}
\caption{The equation in the moving frame holds to numerical precision only when the connection term is retained.}
\label{core:fig:moving-frame}
\end{figure}

\begin{certificate}[Moving-frame correction]\label{core:certificate:005}
The reported residual for the complete equation in the moving frame was $1.0645\times10^{-10}$. Omitting the connection increased the residual to $0.7000$.
\end{certificate}

\subsection{A canonical projector connection}

Let $P_1(t),\ldots,P_m(t)$ be differentiable on an open real time interval, with $P_a(t)^*=P_a(t)$ and
\[
P_aP_b=\delta_{ab}P_a,
\qquad
\sum_aP_a=I.
\]

\begin{theorem}[Canonical packet connection]\label{core:theorem:011}\label{core:thm:canonical}
Define
\[
\boxed{
C(t)
=
\frac12\sum_a[\dot P_a(t),P_a(t)]
=
\sum_a\dot P_a(t)P_a(t).
}
\]
Then
\[
C(t)^*=-C(t)
\]
and
\[
\boxed{
[C(t),P_b(t)]=\dot P_b(t)
}
\]
for every packet $b$.
\end{theorem}

\begin{proof}\phantomsection\label{core:proof:012}
Differentiating $P_a^2=P_a$ gives
\[
\dot P_aP_a+P_a\dot P_a=\dot P_a.
\]
Summing over $a$ and using $\sum_a\dot P_a=0$ yields
\[
\sum_aP_a\dot P_a=-\sum_a\dot P_aP_a,
\]
which proves both displayed formulas for $C$ and its skew-Hermiticity.

For fixed $b$,
\[
CP_b=\dot P_bP_b.
\]
When $a\ne b$, differentiating $P_bP_a=0$ gives
\[
P_b\dot P_a=-\dot P_bP_a.
\]
Hence
\[
P_bC
=\sum_aP_b\dot P_aP_a
=-\dot P_b(I-P_b),
\]
and subtraction gives $[C,P_b]=\dot P_b$.
\phantomsection\label{core:proof:012:end}\end{proof}

The connection generates infinitesimal transport of the moving packet frame. The packet family determines it directly, so its construction does not require a basis inside any packet.

\subsection{Projector curvature}

\begin{classical}[Projected connection]\label{core:classical:001}\label{core:thm:projected}
Let $M$ be a smooth finite-dimensional real manifold without boundary, let $V$ be a finite-dimensional real normed vector space, and let $P:M\to\operatorname{End}_{\bbR}(V)$ be smooth with $P(x)^2=P(x)$. In the quantum application, $V$ is the underlying real space of a finite complex Hilbert space and $P$ is an orthogonal complex-linear projector.

The rank of $P$ is locally constant. On each open rank region, its ranges form a smooth vector bundle. For a smooth range section $s$, so that $Ps=s$, define
\[
\nabla_Xs=P\,ds(X).
\]
This is a connection on smooth range sections. For smooth vector fields $X,Y$, its curvature satisfies
\[
\boxed{
(\nabla_X\nabla_Y-\nabla_Y\nabla_X-\nabla_{[X,Y]})s
=P[dP(X),dP(Y)]P\,s.
}
\]
Equivalently, its endomorphism-valued curvature is
\[
\boxed{F=P(dP\wedge dP)P.}
\]
No global frame or globally constant rank is required, and the vector fields need not commute. The geometric formalization concerns real smooth structures. It does not cover complex analytic manifolds or extensions to a boundary.
\end{classical}

\begin{proof}\phantomsection\label{core:proof:013}
\textit{Range coordinates.} Fix $a\in M$ and put
\[
T_a(x)=P(x)P(a)+(I-P(x))(I-P(a)).
\]
Then $T_a(a)=I$ and $P(x)T_a(x)=T_a(x)P(a)$. Invertibility is open, so on a neighborhood of $a$, $T_a(x)$ restricts to an isomorphism from $\Ran P(a)$ to $\Ran P(x)$. This proves local constancy of rank. The local coordinates and their transitions are
\[
(x,v)\longmapsto (x,T_a(x)^{-1}v),\qquad
u\longmapsto T_b(x)^{-1}T_a(x)u.
\]
Their smoothness follows from smoothness of $P$ and inversion. The transition maps satisfy the cocycle identity by cancellation. These charts construct the smooth range bundles on the open rank regions, including when different connected components have different ranks.

\textit{Connection rules.} Smoothness of $P\,ds(X)$ and $P^2=P$ show that $\nabla_Xs$ is a smooth range section. Differentiation and projection give
\[
\begin{aligned}
\nabla_{X+Y}s&=\nabla_Xs+\nabla_Ys,&
\nabla_{fX}s&=f\nabla_Xs,\\
\nabla_X(s+t)&=\nabla_Xs+\nabla_Xt,&
\nabla_X(fs)&=f\nabla_Xs+(Xf)s
\end{aligned}
\]
for smooth real $f$. The range coordinates above identify ambient smooth range sections with smooth sections of the constructed bundle.

\textit{Curvature.} Differentiating $P^2=P$ gives $P(dP(X))P=0$. Write $D_Xs=ds(X)$. For ambient smooth sections,
\[
D_XD_Ys-D_YD_Xs-D_{[X,Y]}s=0.
\]
Expansion of the projected derivatives therefore leaves
\[
R(X,Y)s=P(dP(X))D_Ys-P(dP(Y))D_Xs.
\]
Since $Ps=s$, differentiation gives $D_Ys=(dP(Y))s+PD_Ys$. The terms with a middle $P$ vanish. Thus
\[
R(X,Y)s=P\bigl(dP(X)dP(Y)-dP(Y)dP(X)\bigr)P\,s.
\]
This expression is alternating and bilinear in the tangent vectors, and it preserves the range. The ordinary chain rule and the transformation law for the Lie bracket make the local calculations agree on chart overlaps, which gives the global formula.
\phantomsection\label{core:proof:013:end}\end{proof}

For an orthogonal rank-one complex projector, use the scalar connection convention $A=-i\langle u,du\rangle$ in a local unit frame. The associated scalar curvature density is
\[
\cF_{\mu\nu}
=-i\Tr\left(P[\partial_\mu P,\partial_\nu P]\right).
\]
On the Bloch sphere,
\[
P(\theta,\phi)
=
\frac12\left(I+\mathbf n(\theta,\phi)\cdot\boldsymbol\sigma\right),
\]
where
\[
\mathbf n=(\sin\theta\cos\phi,\sin\theta\sin\phi,\cos\theta).
\]
The Pauli multiplication law gives
\[
[\partial_\theta P,\partial_\phi P]
=\frac{i}{2}(\partial_\theta\mathbf n\mathbin{\times}\partial_\phi\mathbf n)\cdot\boldsymbol\sigma.
\]
Using $\Tr(P\,\mathbf a\cdot\boldsymbol\sigma)=\mathbf n\cdot\mathbf a$ therefore yields
\[
\boxed{\cF_{\theta\phi}
=\frac12\mathbf n\cdot(\partial_\theta\mathbf n\mathbin{\times}\partial_\phi\mathbf n)
=\frac12\sin\theta.}
\]
The scalar sign changes if the opposite convention for $A$ is chosen. This explicit calculation is not a separate entry among the 19 items in the Lean register.

\begin{figure}[htbp]
\centering
\includegraphics[width=0.76\textwidth]{09_bloch_curvature.pdf}
\caption{The curvature of the Bloch projector, evaluated by finite differences, is compared with the exact density $\frac12\sin\theta$.}
\label{core:fig:bloch}
\end{figure}

\begin{certificate}[Geometric packet identities]\label{core:certificate:006}
The canonical connection reconstructed every packet derivative with residual $0$. The maximum residual for the Bloch curvature was $2.0803\times10^{-11}$.
\end{certificate}

\section{Measurements and quantum channels}\label{core:section:007}

\subsection{Packet measurements}

Let $\rho$ be a density matrix and let $\{E_k\}$ be a POVM:
\[
E_k\ge0,
\qquad
\sum_kE_k=I.
\]
The outcome probabilities are
\[
p_k=\Tr(E_k\rho).
\]

\begin{theorem}[Block-diagonal measurement invariance]\label{core:theorem:012}\label{core:thm:measurement}
Let $\{P_\alpha\}$ be a complete packet family, and suppose every measurement effect is packet diagonal:
\[
E_k=\sum_\alpha P_\alpha E_kP_\alpha.
\]
Define the packet-dephased state
\[
\cD_P(\rho)=\sum_\alpha P_\alpha\rho P_\alpha.
\]
Then
\[
\boxed{
\Tr(E_k\rho)
=
\Tr(E_k\cD_P(\rho))
}
\]
for every outcome $k$.
\end{theorem}

\begin{proof}\phantomsection\label{core:proof:014}
Insert the packet decomposition of $E_k$ and use cyclicity of the trace:
\[
\Tr(E_k\rho)
=
\sum_\alpha\Tr(P_\alpha E_kP_\alpha\rho)
=
\sum_\alpha\Tr(E_kP_\alpha\rho P_\alpha).
\]
\phantomsection\label{core:proof:014:end}\end{proof}

The theorem identifies which coherences a declared measurement can detect. Observables that are diagonal in the packet decomposition cannot detect coherence between distinct packets. Coherence within a packet remains available to such observables.

\subsection{Kraus transfer}

A quantum channel can transfer both populations and coherences between packet addresses. A Kraus representation describes this transfer by summing the contributions of its individual branches. The normalization of those branches ensures that the total trace is preserved.

\begin{classical}[Finite Kraus channel]\label{core:classical:002}\label{core:thm:kraus}
Let $K_1,\ldots,K_r$ satisfy
\[
\sum_aK_a^*K_a=I.
\]
Then
\[
\boxed{
\Phi(\rho)=\sum_aK_a\rho K_a^*
}
\]
is completely positive and trace preserving.
\end{classical}

\begin{proof}\phantomsection\label{core:proof:015}
Trace preservation follows from
\[
\Tr\Phi(\rho)
=\Tr\left(\rho\sum_aK_a^*K_a\right)
=\Tr\rho.
\]
For any auxiliary dimension and positive matrix $X$, each term
\[
(I\otimes K_a)X(I\otimes K_a)^*
\]
is positive, so their sum is positive. This is complete positivity.
\phantomsection\label{core:proof:015:end}\end{proof}

Packet superoperator blocks are defined by
\[
\Phi_{\alpha\beta;\gamma\delta}(X)
=
P_\alpha
\Phi(P_\gamma XP_\delta)
P_\beta.
\]
Summing all four packet indices reconstructs $\Phi(X)$ exactly.

For amplitude damping with a real parameter $0\le\gamma\le1$,
\[
K_0=
\begin{pmatrix}
1&0\\0&\sqrt{1-\gamma}
\end{pmatrix},
\qquad
K_1=
\begin{pmatrix}
0&\sqrt\gamma\\0&0
\end{pmatrix}.
\]
The excited state becomes
\[
\Phi(\ket1\bra1)
=
\gamma\ket0\bra0+(1-\gamma)\ket1\bra1.
\]

\begin{figure}[htbp]
\centering
\includegraphics[width=0.56\textwidth]{16_amplitude_damping_channel.pdf}
\caption{Amplitude damping transfers population from the excited packet to the ground packet while preserving trace and positivity.}
\label{core:fig:damping}
\end{figure}

\begin{certificate}[Channel reconstruction]\label{core:certificate:007}
For $\gamma=0.23$, the residuals for Kraus completeness and the output trace were $0$. The computed output from $\ket1\bra1$ was $\diag(0.23,0.77)$. In the Bell state test, packet dephasing left every measurement probability in the computational basis unchanged.
\end{certificate}

\section{Quantum error correction as obstruction-packet annihilation}\label{core:section:008}

Let $P=P^*=P^2$ project onto a code space $\cC\subset\cH$, and let $E_1,\ldots,E_r$ be a finite family of complex error matrices. Neither a nonzero code nor channel normalization is assumed unless stated. The operators
\[
P E_a^*E_bP
\]
are the obstruction packets on the code. If they retain information about the logical state beyond one scalar multiple of $P$, the environment can distinguish logical states and exact correction is impossible.

We call a recovery complete when it is a Kraus map
\[
\cR(X)=\sum_jK_jXK_j^*,\qquad \sum_jK_j^*K_j=I.
\]
The family is \emph{exactly correctable} when one such recovery works for every coherent error combination $A_w=\sum_aw_aE_a$: for every $w$ there is a real $t_w\ge0$, independent of the logical input, such that
\[
PXP=X\quad\Longrightarrow\quad
\cR(A_wXA_w^*)=t_wX.
\]
The coefficient is an unnormalized weight. It is not necessarily a probability for an arbitrary matrix $A_w$.

\begin{classical}[Finite Knill--Laflamme criterion]\label{core:classical:003}\label{core:thm:kl}
The finite error family is exactly correctable in the sense just defined if and only if there is a positive semidefinite complex matrix $C=(c_{ab})$ such that
\[
\boxed{PE_a^*E_bP=c_{ab}P\qquad\text{for every }a,b.}
\]
If $\sum_aE_a^*E_a=I$, the equivalent normalized-channel formulation is the existence of a complete recovery satisfying
\[
\cR\!\left(\sum_aE_aXE_a^*\right)=X
\quad\text{for every }X\text{ with }PXP=X.
\]
The zero code and an empty error family are included. An empty family cannot define a channel that preserves trace on a nonzero ambient space.
\end{classical}

\begin{proof}\phantomsection\label{core:proof:016}
\textit{Necessity.} Fix an error $E_a$ and let $L_j=K_jE_aP$. Perfect recovery on rank-one code operators gives
\[
\sum_j\ket{L_j\psi}\bra{L_j\psi}
=t_a\ket\psi\bra\psi\qquad(\psi\in\cC).
\]
For every vector orthogonal to $\psi$, the quadratic form on the left is a sum of nonnegative squared absolute values and equals zero. Each $L_j\psi$ therefore lies in the span of $\psi$. A linear operator that preserves every line in $\cC$ is scalar there: apply it to two independent vectors and their sum. The one-dimensional case is immediate, and the claim is vacuous for the zero code. Consequently
\[
K_jE_aP=d_{ja}P.
\]
Recovery completeness now gives
\[
\begin{aligned}
PE_a^*E_bP
&=\sum_j (K_jE_aP)^*(K_jE_bP)\\
&=\left(\sum_j\overline{d_{ja}}d_{jb}\right)P.
\end{aligned}
\]
The coefficients form a Gram matrix, hence are positive semidefinite. For a normalized channel, apply the same rank-one argument to all composite branches $K_jE_aP$ at once to obtain the necessary condition.

\textit{Sufficiency.} Diagonalize $C=U\operatorname{diag}(\lambda_r)U^*$ with $U$ unitary and $\lambda_r\ge0$. Define compressed transformed errors
\[
F_r=\sum_aU_{ar}E_aP.
\]
Then $F_rP=F_r$ and $F_r^*F_s=\lambda_r\delta_{rs}P$. If $\lambda_r=0$, $F_r^*F_r=0$ forces $F_r=0$. For positive eigenvalues only, set $V_r=F_r/\sqrt{\lambda_r}$. These satisfy
\[
V_rP=V_r,\qquad V_r^*V_s=\delta_{rs}P.
\]
Let $S=\sum_{\lambda_r>0}V_rV_r^*$. Orthogonality makes $S$ an orthogonal projector. Use the recovery branches
\[
K_r=V_r^*,\qquad K_{\mathrm{unused}}=I-S.
\]
They are complete on the whole ambient space because
\[
\sum_rK_r^*K_r+K_{\mathrm{unused}}^*K_{\mathrm{unused}}
=S+(I-S)^2=I.
\]
The inverse unitary change of coordinates reconstructs each $E_aP$ from the $F_r$. Every recovery branch of a coherent combination is therefore scalar on the code, while the branch for the unused space vanishes. For all supported matrices, the resulting recovery is
\[
\cR(A_wXA_w^*)=(w^*Cw)X.
\]
This establishes the arbitrary-family statement without assuming normalization. Under $\sum_aE_a^*E_a=I$ and $P\ne0$, summing the diagonal KL equations gives $\Tr C=1$, so the normalized-channel output is $X$. When $P=0$, every supported $X$ is zero and the conclusion is automatic. The recovery is trace preserving by completeness and completely positive at every finite ancillary dimension by the Kraus calculation in Classical Input~\ref{core:thm:kraus}.
\phantomsection\label{core:proof:016:end}\end{proof}

For the three-qubit bit-flip code,
\[
\ket{0_L}=\ket{000},
\qquad
\ket{1_L}=\ket{111},
\]
and the error family is
\[
I,
\quad X\otimes I\otimes I,
\quad I\otimes X\otimes I,
\quad I\otimes I\otimes X.
\]
The four error images of the code are pairwise orthogonal two-dimensional syndrome packets whose projectors sum to the identity on the eight-dimensional Hilbert space.

\begin{figure}[htbp]
\centering
\includegraphics[width=0.88\textwidth]{19_error_correction_syndromes.pdf}
\caption{The four orthogonal syndrome packets of the three-qubit bit-flip code. Each packet identifies one correctable error without resolving the encoded logical state.}
\label{core:fig:error-correction}
\end{figure}

\begin{certificate}[Exact bit-flip correction]\label{core:certificate:008}
The residuals for the Knill--Laflamme condition and completeness of the syndrome packets were both $0$. A mixed channel of single bit flip errors, applied to a nontrivial logical superposition, was reversed with a residual of $1.0385\times10^{-16}$ in the density matrix.
\end{certificate}

\section{Exact effective dynamics}\label{core:section:009}

\subsection{Schur--Feshbach reduction}

Let $P$ be a retained packet projector and $Q=I-P$. In packet coordinates,
\[
H=
\begin{pmatrix}
H_{PP}&H_{PQ}\\
H_{QP}&H_{QQ}
\end{pmatrix}.
\]

\begin{theorem}[Exact Feshbach resolvent]\label{core:theorem:013}\label{core:thm:feshbach}
Assume that both $zI_Q-H_{QQ}$ and $zI-H$ are invertible. Define
\[
\boxed{
H_{\mathrm{eff}}(z)
=
H_{PP}
+H_{PQ}(z-H_{QQ})^{-1}H_{QP}.
}
\]
Then
\[
\boxed{
P(z-H)^{-1}P
=
\left(z-H_{\mathrm{eff}}(z)\right)^{-1}
}
\]
where the left side is restricted to $P\cH$, the right side is an inverse on that retained space, and both are written in the same orthonormal basis. Here $z-H_{QQ}$ and $z-H_{\mathrm{eff}}$ abbreviate the appropriately sized scalar matrices minus the indicated blocks.
\end{theorem}

\begin{proof}\phantomsection\label{core:proof:017}
The block in the lower right corner of $z-H$ is $z-H_{QQ}$. Gaussian block elimination produces the Schur complement
\[
z-H_{PP}-H_{PQ}(z-H_{QQ})^{-1}H_{QP}.
\]
The assumed invertibility of this block and of the full matrix implies that the Schur complement is invertible. The block in the upper left corner of the full inverse is the inverse of that Schur complement, which is exactly the displayed operator.
\phantomsection\label{core:proof:017:end}\end{proof}

The second premise cannot be dropped: at $z=0$ and $H=\operatorname{diag}(0,-1)$, the eliminated block is $1$ but $zI-H$ is singular. This counterexample is also included in the Lean module.

The self-energy
\[
\Sigma_P(z)
=H_{PQ}(z-H_{QQ})^{-1}H_{QP}
\]
contains the complete influence of the eliminated packet. Bare truncation drops $\Sigma_P(z)$ and is therefore exact only when the coupling between packets vanishes or the self-energy cancels for another stated reason.

\begin{figure}[htbp]
\centering
\begin{minipage}[t]{0.48\textwidth}
\centering
\includegraphics[width=\textwidth]{06_feshbach_vs_truncation.pdf}
\captionof{figure}{Exact Feshbach reduction versus bare truncation.}
\label{core:fig:feshbach}
\end{minipage}\hfill
\begin{minipage}[t]{0.48\textwidth}
\centering
\includegraphics[width=\textwidth]{07_packet_green_blocks.pdf}
\captionof{figure}{Norms of packet Green-function blocks.}
\label{core:fig:green-blocks}
\end{minipage}
\end{figure}

\begin{certificate}[Effective resolvent]\label{core:certificate:009}
In the reported test with a random Hermitian matrix of dimension six, the exact Feshbach residual was $3.2291\times10^{-17}$. The corresponding residual for bare truncation was $1.3091\times10^{-2}$.
\end{certificate}

\subsection{The hidden packet as memory}

For a time-independent finite Hamiltonian $H$ and fixed projectors $P,Q$, write $\psi=x+y$, with $x=P\psi$ and $y=Q\psi$. Then
\[
\begin{aligned}
i\dot x&=H_{PP}x+H_{PQ}y,\\
i\dot y&=H_{QP}x+H_{QQ}y.
\end{aligned}
\]
Solving the second equation gives
\[
y(t)
=e^{-itH_{QQ}}y(0)
-i\int_0^t e^{-i(t-s)H_{QQ}}H_{QP}x(s)\,ds.
\]
Substitution yields the exact Volterra equation
\[
\boxed{
\begin{aligned}
i\dot x(t)
={}&H_{PP}x(t)
+H_{PQ}e^{-itH_{QQ}}y(0)\\
&-i\int_0^t
K(t-s)x(s)\,ds,
\end{aligned}
}
\]
with memory kernel
\[
\boxed{
K(t)=H_{PQ}e^{-itH_{QQ}}H_{QP}.
}
\]
The influence of the eliminated degrees of freedom is therefore retained in a self-energy that depends on frequency, or equivalently in an exact memory term. For forward evolution the integral is causal. When the Hamiltonian depends on time, the corresponding propagator between two times replaces the stationary exponential used here. This Volterra derivation is not a separate entry in the finite Lean register.

\section{Finite quantum geometry}\label{core:section:010}

\subsection{A projector-resolved geometry}

At resolution $\ell$, define
\[
\boxed{
\mathfrak Q_\ell
=
\left(
\cA_\ell,
\cH_\ell,
D_\ell,
\{P_{\ell,\alpha}\},
E_\ell,
C_\ell,
\Pi_\ell^{\mathrm{phys}},
T_\ell
\right).
}
\]
In order, the entries specify the observable algebra, Hilbert space, geometric or Dirac-type operator, packet projectors, embedding into the fine space, reconstruction onto the coarse space, physical constraint projector, and transfer operator. The maps have types $E_\ell:\cH_\ell\to\cH_{\ell+1}$ and $C_\ell:\cH_{\ell+1}\to\cH_\ell$. The left inverse relation $C_\ell E_\ell=I$ alone does not imply $C_\ell=E_\ell^*$ or make the embedding isometric.

A coherent refinement should satisfy
\[
C_\ell E_\ell=I
\]
and
\[
\boxed{
P_{\ell,\alpha}
=
C_\ell
\left(
\sum_{\beta\succ\alpha}P_{\ell+1,\beta}
\right)
E_\ell.
}
\]
For an operator $D$, define the refinement defect
\[
\Delta_\ell(D)
=C_\ell D_{\ell+1}E_\ell-D_\ell.
\]
Exact resolution independence is $\Delta_\ell(D)=0$. When exact equality is too strong, the correct continuum problem is to prove convergence of the resolvents and uniform control of the effective self-energies.

\begin{figure}[htbp]
\centering
\includegraphics[width=0.84\textwidth]{17_quantum_spacetime_tower.pdf}
\caption{The diagram organizes a conjectural nested packet tower for quantum geometry. It does not provide a numerical reconstruction of physical spacetime.}
\label{core:fig:spacetime-tower}
\end{figure}
\FloatBarrier

\subsection{Euclidean transfer}

\begin{theorem}[Finite Euclidean transfer]\label{core:theorem:014}\label{core:thm:euclidean}
Let $H$ be Hermitian on a nonzero finite-dimensional Hilbert space. For all real $\tau,\sigma\ge0$,
\[
T(\tau)=e^{-\tau(H-E_0I)},
\qquad
E_0=\min\Spec H,
\]
is positive definite and satisfies
\[
\boxed{
T(\tau+\sigma)=T(\tau)T(\sigma).
}
\]
Its spectrum is contained in $(0,1]$.
\end{theorem}

\begin{proof}\phantomsection\label{core:proof:018}
Diagonalize $H=\sum_jE_jP_j$. Then
\[
T(\tau)=\sum_j e^{-\tau(E_j-E_0)}P_j.
\]
Every eigenvalue is strictly positive and at most one. Multiplication adds the exponents, proving the semigroup law and $T(0)=I$. Positivity here is operator positivity, not an assertion that every matrix entry is nonnegative.
\phantomsection\label{core:proof:018:end}\end{proof}

The paper tests a controlled volume-state model
\[
H_{nn}=2\kappa+V(n),
\qquad
H_{n,n+1}=H_{n+1,n}=-\kappa,
\]
with $\kappa>0$ and a real confining potential. This is a finite transfer-matrix test only. Operator positivity alone is not reflection positivity of a Euclidean field measure. It is not derived from a gravitational path integral and is not evidence for a unique microscopic spacetime dynamics.

\begin{figure}[htbp]
\centering
\begin{minipage}[t]{0.48\textwidth}
\centering
\includegraphics[width=\textwidth]{13_euclidean_transfer_spectrum.pdf}
\captionof{figure}{Positive spectrum of the Euclidean transfer operator.}
\label{core:fig:euclidean-spectrum}
\end{minipage}\hfill
\begin{minipage}[t]{0.48\textwidth}
\centering
\includegraphics[width=\textwidth]{14_causal_volume_ground_state.pdf}
\captionof{figure}{Ground-state volume packet in the controlled toy model.}
\label{core:fig:volume-ground}
\end{minipage}
\end{figure}

\begin{figure}[htbp]
\centering
\includegraphics[width=0.68\textwidth]{15_causal_volume_transfer.pdf}
\caption{The figure shows the finite Euclidean kernel for volume transfer. Positivity and semigroup composition hold exactly in the finite model. Its gravitational interpretation remains provisional.}
\label{core:fig:volume-transfer}
\end{figure}

\begin{certificate}[Finite transfer model]\label{core:certificate:010}
The minimum transfer eigenvalue was $0.0767054$, the semigroup residual was $0$, and reconstruction from four volume-residue packets had residual $0$.
\end{certificate}

\section{Spectral dimension and continuum diagnostics}\label{core:section:011}

Let $L\ge0$ be a finite Laplacian on $N\ge1$ states, with eigenvalues $\lambda_j$ and at least one zero eigenvalue. For diffusion time $\tau>0$, define the normalized heat trace
\[
K(\tau)=\frac1N\sum_j e^{-\tau\lambda_j}
\]
and the scale-dependent spectral dimension
\[
\boxed{
d_s(\tau)
=-2\frac{d\log K(\tau)}{d\log\tau}.
}
\]
For a finite system, $d_s(\tau)$ is a diagnostic, not a topological invariant. Indeed, finite differentiation gives
\[
d_s(\tau)=2\tau\,
\frac{\sum_j\lambda_j e^{-\tau\lambda_j}}{\sum_j e^{-\tau\lambda_j}}.
\]
It tends to zero as $\tau\downarrow0$ and, under the stated assumption of a zero mode, as $\tau\to\infty$. If a general positive operator has no zero mode, its behavior at large times is instead asymptotic to $2\tau\lambda_{\min}$. The conclusion for operators with a zero mode therefore does not apply to that case. Inferring a geometric dimension requires a stable intermediate window and a controlled refinement limit.

\subsection{The circle dispersion bound}

For a periodic grid with spacing $h>0$ and physical wave number $k\in\bbR$,
\[
\lambda_{k,h}
=
\frac4{h^2}\sin^2\left(\frac{kh}{2}\right).
\]

\begin{theorem}[Low-mode second-order bound]\label{core:theorem:015}\label{core:thm:lowmode}
For every real $k$ and every nonzero real $h$, define $\lambda_{k,h}=4h^{-2}\sin^2(kh/2)$. Then
\[
\boxed{
0\le k^2-\lambda_{k,h}
\le
\frac{k^4h^2}{12}.
}
\]
\end{theorem}

\begin{proof}\phantomsection\label{core:proof:019}
Put $x=kh/2$. The inequalities $|\sin x|\le|x|$ and $|x-\sin x|\le |x|^3/6$, which hold for every real argument, give
\[
0\le x^2-\sin^2x
=|x-\sin x|\,|x+\sin x|
\le\frac{|x|^3}{6}\,2|x|
=\frac{x^4}{3}.
\]
Multiplying by $4/h^2$ proves both bounds. The estimate holds for every frequency. Its interpretation as convergence of second order keeps $k$ fixed while $h\to0$, so it does not assert a uniformly small relative error across all grid modes.
\phantomsection\label{core:proof:019:end}\end{proof}

\begin{figure}[htbp]
\centering
\includegraphics[width=0.76\textwidth]{11_circle_low_mode_convergence.pdf}
\caption{The first nonzero circle mode converges at second order and remains below the proved dispersion bound.}
\label{core:fig:circle-convergence}
\end{figure}

\begin{certificate}[Continuum dispersion]\label{core:certificate:011}
The reported convergence order was $1.99936$. For every tested mode, the error was below the bound, so none of those tests violated the theorem.
\end{certificate}

\subsection{One- and two-dimensional plateaus}

At intermediate diffusion times, a finite cycle and a finite torus can exhibit plateaus near the dimensions of their underlying geometries. The following figure and numerical certificate show the reported behavior for these two examples. The plateaus occur between the small and large time regimes in which the spectral dimension returns to zero.

\begin{figure}[htbp]
\centering
\includegraphics[width=0.80\textwidth]{12_spectral_dimension.pdf}
\caption{The curves show intermediate windows of spectral dimension for a finite cycle and a finite two-dimensional torus. Departures in the ultraviolet and infrared reflect the finite size of the models rather than additional dimensions.}
\label{core:fig:spectral-dimension}
\end{figure}

\begin{certificate}[Spectral dimension]\label{core:certificate:012}
The median intermediate plateaus were
\[
d_s^{(1)}=1.00062,
\qquad
d_s^{(2)}=2.03781.
\]
These tests recover the expected dimensions of the chosen finite geometries. They do not select a physical theory of quantum spacetime.
\end{certificate}

\section{Constraint packets and the quantum-spacetime boundary}\label{core:section:012}

At one finite level, let $\widehat C_A$ be a finite family of matrix constraints, with fixed scalar structure coefficients, and let $\Pi^{\mathrm{phys}}$ be the orthogonal projector onto their exact common kernel. If the classical constraint algebra has structure coefficients $f_{AB}^{\ \ C}$, define the anomaly operator
\[
\boxed{
\mathfrak A_{AB}
=
[\widehat C_A,\widehat C_B]
-i\sum_C f_{AB}^{\ \ C}\widehat C_C.
}
\]
Its physical packet is
\[
\Pi^{\mathrm{phys}}\mathfrak A_{AB}\Pi^{\mathrm{phys}}.
\]
In this exact finite setting, $\widehat C_C\Pi^{\mathrm{phys}}=0$ for every $C$. Associativity gives
\[
\mathfrak A_{AB}\Pi^{\mathrm{phys}}=0,
\qquad
\Pi^{\mathrm{phys}}\mathfrak A_{AB}\Pi^{\mathrm{phys}}=0
\]
automatically. This compression therefore cannot certify operator closure away from the physical range. A nonzero measured compression contradicts the exact-kernel assumptions or indicates an approximation or implementation error. For example, $C_1=\operatorname{diag}(1,0,0)$ and $C_2=E_{12}+E_{21}$ annihilate the same third coordinate, while their commutator $E_{12}-E_{21}$ is nonzero. With zero proposed structure coefficients, the anomaly is nonzero but its physical compression vanishes.

The closure target must be stated at the operator level, not inferred from that automatic compression. Unbounded constraints additionally require a specified common invariant domain. Neither such a domain theory nor gravitational structure functions are supplied by the finite check.

\begin{conjecture}[Anomaly-free refinement]\label{core:conjecture:001}
For a specified nontrivial quantum-geometric constraint model, there exists a coherent packet tower with nonzero physical spaces, compatible physical projectors and exact operator closure
\[
\mathfrak A^{(\ell)}_{AB}=0
\]
at every finite level. The compatible physical compression then vanishes as a consequence, not as a substitute for closure.
\end{conjecture}
The substantive research target is the operator condition. The earlier condition stated only in terms of compression was too weak. A choice of model, dynamics, and refinement maps is still needed before this program becomes a fully specified physical existence problem.

\begin{conjecture}[Uniform Feshbach control]\label{core:conjecture:002}
For the same specified tower, the self-energies from exact elimination of fine packets admit bounds uniform in cutoff and volume on a common spectral domain. Together with consistent identifications of the Hilbert spaces and a proved resolvent convergence criterion, they yield a nontrivial continuum limit in the strong resolvent sense. Uniform boundedness alone is not asserted to imply convergence.
\end{conjecture}

\begin{conjecture}[Lorentzian reconstruction]\label{core:conjecture:003}
The Euclidean packet theory satisfies a positivity or reconstruction principle strong enough to produce a Lorentzian Hilbert space with causal propagation and positive physical norm.
\end{conjecture}

\begin{conjecture}[Einstein infrared limit]\label{core:conjecture:004}
There exist semiclassical packet states $\Psi_{g,\ell}$ such that the effective action converges to
\[
\frac1{16\pi G}
\int\sqrt{-g}\,(R-2\Lambda)\,d^4x
\]
with controlled higher-order corrections.
\end{conjecture}

\begin{conjecture}[Massless spin-two sector]\label{core:conjecture:005}
The continuum two-point function possesses a positive transverse-traceless massless spin-two pole, while all unphysical packet sectors are removed by the constraints.
\end{conjecture}

The exact finite identities establish none of these physical conclusions, and no formal theorem of logical independence is claimed. Exact packet reconstruction provides a kinematic and computational advantage. It does not determine the physical action, prove existence of the continuum theory, or guarantee the limit of general relativity.

\begin{figure}[htbp]
\centering
\includegraphics[width=0.86\textwidth]{18_proof_boundary.pdf}
\caption{The proof boundary. Finite packet, circuit, leakage, Feshbach, connection, and spectral identities are established here. Constraint closure and the physical continuum claims remain conjectural.}
\label{core:fig:proof-boundary}
\end{figure}

\section{Numerical certificates}\label{core:section:013}

The numerical study reports CPU experiments and a suite of 22 tests. The values below come from the accompanying numerical record. They are separate from the 323 supporting declarations in the formal companion. A displayed floating-point residual of zero describes the reported computation and does not independently prove equality in exact arithmetic.

\begin{table}[htbp]
\centering
\caption{Reported finite numerical checks, retained from the original manuscript and not rerun for this edition.}
\label{core:tab:certificates}
\begin{tabular}{L{0.64\textwidth}r}
\toprule
Certificate & Value\\
\midrule
Bell-state residual & $0$\\
Three-qubit seed-orbit rank & $8/8$\\
Parity off-block norm & $0$\\
Circuit path reconstruction residual & $0$\\
Exact-symmetry maximum leakage & $0$\\
Maximum leakage-bound violation & $0$\\
Feshbach projected-resolvent residual & $3.2291\times10^{-17}$\\
Bare-truncation projected-resolvent residual & $1.3091\times10^{-2}$\\
Moving-frame residual with connection & $1.0645\times10^{-10}$\\
Moving-frame residual without connection & $7.0000\times10^{-1}$\\
Canonical packet-connection residual & $0$\\
Bloch-curvature residual & $2.0803\times10^{-11}$\\
Kraus completeness residual & $0$\\
Output trace residual & $0$\\
Knill--Laflamme residual & $0$\\
Bit-flip recovery residual & $1.0385\times10^{-16}$\\
Tensor parity-refinement residual & $0$\\
Circle low-mode convergence order & $1.99936$\\
One-dimensional spectral plateau & $1.00062$\\
Two-dimensional spectral plateau & $2.03781$\\
Minimum Euclidean-transfer eigenvalue & $0.0767054$\\
Transfer semigroup residual & $0$\\
Volume-packet reconstruction residual & $0$\\
\bottomrule
\end{tabular}
\end{table}

The original reported count is
\[
\boxed{22/22\text{ historical numerical tests, not rerun here}.}
\]
The 19 original figures are retained from the recovered numerical companion. The separate visual guide supplies a reading route through these constructions and the later strengthened results.

\setcounter{section}{14}
\section{Conclusion}\label{core:section:014}

The quantum realization of the seed calculus begins with a primitive state and a declared operator grammar. Symmetry projectors resolve the state and every operator into exact packets. In these coordinates, circuit products become sums over packet paths, and Hamiltonian symmetry appears as a diagonal block structure. When that symmetry is broken, the resulting transfer is measured by the leakage estimates.

Moving packets carry a canonical connection and curvature. Channels and measurements act on the corresponding operator packets, while the conditions for error correction become statements about compression to the code. Exact elimination retains the influence of hidden sectors through self-energy and memory.

The finite theory therefore provides a coherent language for
\[
\boxed{
\begin{gathered}
\text{seed reachability}
+
\text{symmetry resolution}
+
\text{circuit paths}\\
+
\text{effective dynamics}
+
\text{geometric transport}
+
\text{constraint diagnostics}
\end{gathered}
}
\]
A physical continuum theorem remains open. The next required results are anomaly-free refinement, uniform control of effective operators, Lorentzian reconstruction, the Einstein infrared limit, and a positive massless spin-two sector. Each requires its own hypotheses and proof. Agreement in a finite numerical model cannot replace those arguments.


